% TOP_PROBLEM: {"id": 3218, "title": "Bouchet's 6-flow conjecture", "category_id": 3, "category": {"id": 3, "name": "graph_theory", "display_name": "Graph Theory", "description": "Problems involving graphs, networks, and their properties.", "slug": "graph-theory", "order_index": 3, "created_at": "2026-07-31T15:26:25.671Z"}, "status": "open", "problem_number": "OPG-131", "set_id": 12}
% FIRST_POSTED: 2026-10-01
% ENTRY_KIND: statement_only
% UnsolvedMath ID 3218; source catalogue code OPG-131
% !TeX program = lualatex
% Definition and sources only; this file is not an LLM solution attempt.
\documentclass[11pt,a4paper]{article}
\usepackage[margin=25mm]{geometry}
\usepackage{fontspec}
\setmainfont{lmroman10-regular.otf}[BoldFont=lmroman10-bold.otf,
 ItalicFont=lmroman10-italic.otf,BoldItalicFont=lmroman10-bolditalic.otf]
\usepackage{amsmath,amssymb,mathtools}
\usepackage{unicode-math}
\setmathfont{latinmodern-math.otf}
\AtBeginDocument{\let\setminus\smallsetminus}
\usepackage{xurl,hyperref}
\hypersetup{colorlinks=true,urlcolor=blue}
\setcounter{secnumdepth}{0}
\setlength{\parindent}{0pt}
\setlength{\parskip}{5pt}
\setlength{\emergencystretch}{3em}
\begin{document}
\section*{Bouchet's 6-flow conjecture}\label{3218}
{\small UnsolvedMath ID 3218 · OPG-131 · Graph Theory · OpenGarden}
\subsection{Definitions and mathematical statement}
A finite bidirected multigraph consists
of an undirected multigraph together
with a sign on each incidence,
or half-edge, at a vertex. Use
sign $+1$ for an end pointing
into its vertex and $-1$ for
an end pointing out. The two
ends of an edge are chosen
independently: both may point in,
both out, or in opposite directions.
A loop has two incidences at
the same vertex, each counted
separately in the sums below.

For an integer $q\geq2$, an integer
$q$-flow is a map $\phi:E\to\mathbb Z$
with $|\phi(e)|<q$ and
\[
 \sum_{h\text{ incident with }v}
       \sigma(h)\phi(e(h))=0
 \qquad\text{for every vertex }v,
\]
where $e(h)$ is the edge of
the half-edge $h$. The flow
is nowhere-zero if $\phi(e)\neq0$
for every edge. A bidirected
graph is flow-admissible if it
has a nowhere-zero integer
$q$-flow for some integer $q$.

Bouchet's conjecture states that
every flow-admissible finite
bidirected graph has a nowhere-zero
integer $6$-flow, with values
in $\{-5,-4,-3,-2,-1,1,2,3,4,5\}$.
The bidirected incidence signs
are fixed as part of the
input and must be respected
by the smaller flow.

The premise is existence of
an integer flow of some size,
not just absence of bridges
in the underlying undirected
graph. Unlike ordinary orientations,
an edge can contribute with
the same sign at both endpoints;
the half-edge definition includes
this distinction explicitly.
\subsection{Short English statement}
Whenever a bidirected graph admits any nonzero integer flow, can it admit one with every edge magnitude at most five?
\subsection{Sources}
\begin{itemize}
\item[S1] ULAM.AI and the UnsolvedMath Contributors, supplied problems.json, record 3218 (OPG-131), OpenGarden collection. Catalogue identity and source transcription; CC BY 4.0.
\url{https://huggingface.co/datasets/ulamai/UnsolvedMath}.
\item[S2] Open Problem Garden, Bouchet's 6-flow conjecture. Original problem and discussion; the mathematical formulation below is expanded from the supplied source record.
\url{http://www.openproblemgarden.org/op/bouchets_6_flow_conjecture}.
\end{itemize}
\end{document}
